\chapter{Arithmetic Functions on Average}

\begin{center}\itshape
Die ganzen Zahlen hat der liebe Gott gemacht, alles andere ist Menschenwerk.\\
(God made the integers; all else is the work of man.)\\
\upshape --- Leopold Kronecker, as reported by Heinrich Weber (1893)
\end{center}

Ask a hundred people their income and the answers are all over the place; average them and you get a stable, sensible number,
which almost none of them actually earns. Integers behave the same way. The number of divisors $\tau(n)$ jumps around wildly: $\tau(p)=2$ for every prime, while $\tau(720720)=240$. Euler's function
$\varphi(n)$, the sum of divisors $\sigma(n)$, the indicator $\mu^2$ of the squarefree integers behave just as erratically. Yet
their averages are perfectly regular. Volume~2 (Theorem~4.7) showed that $\tau(n)$ is $\log n$ on average, with Dirichlet's hyperbola
method; Volume~2, Exercise~4.8 and Volume~5, Exercise~12.2 found the densities $\frac6{\pi^2}$ of coprime pairs and of squarefree
integers.

This chapter turns these examples into a method, and the method has a moral worth keeping in mind for the whole volume: \emph{on average,
divisibility by different primes behaves like a set of independent coins.} One theorem (Wintner's) computes the mean value of every function of
the form $\mathbf1*h$ with $h$ small, and explains why the answer is an Euler product, as if divisibility by different primes were
independent events. Then we ask for error terms, iterate the hyperbola method to the three-fold divisor function, and meet the
first of the great open problems of the subject, the size of the error in Dirichlet's divisor problem: a question about counting
lattice points under a hyperbola, which anyone can state on a napkin, and which has been open since 1849.

\begin{think}
The average of $\tau(n)$ over $n\le10^6$ is $\log10^6+2\gamma-1\approx13.97$ (Volume~2, Theorem~4.7). Guess the median of $\tau(n)$ for
$n\le10^6$, and the proportion of $n\le10^6$ with at least $14$ divisors. Is a ``typical'' integer near $10^6$ average?
\end{think}

\section{Mean values}

We keep the notation of Volume~2, Chapter~4: $f*g$ is Dirichlet convolution, $(f*g)(n)=\sum_{dm=n}f(d)g(m)$; $\mathbf1$ is the constant
function $1$, $\operatorname{id}(n)=n$, $\varepsilon$ is the unit ($\varepsilon(1)=1$ and $\varepsilon(n)=0$ for $n>1$), and $\mu*\mathbf1=\varepsilon$. Thus
$\tau=\mathbf1*\mathbf1$, $\sigma=\operatorname{id}*\mathbf1$ and $\varphi=\mu*\operatorname{id}$. An arithmetic function $f$ has \emph{mean value} $m(f)$ if
$\frac1x\sum_{n\le x}f(n)\to m(f)$ as $x\to\infty$.

\begin{lemma}\label{lem:convsum}
Let $g,h$ be arithmetic functions and $G(y)=\sum_{m\le y}g(m)$. Then for $x\ge1$,
\[
\sum_{n\le x}(h*g)(n)=\sum_{d\le x}h(d)\,G\Bigl(\frac xd\Bigr).
\]
\end{lemma}
\begin{proof}
Both sides equal $\sum h(d)g(m)$ over the pairs $(d,m)$ of positive integers with $dm\le x$: on the left grouped by $n=dm$, on the right
by $d$, since $dm\le x$ means $m\le\frac xd$.
\end{proof}

With $g=\mathbf1$, $G(y)=\lfloor y\rfloor$, and $\lfloor\frac xd\rfloor=\frac xd-\{\frac xd\}$. If $h$ is small, the main term $x\sum_d\frac{h(d)}d$ wins.

\begin{theorem}[Wintner, 1943]\label{thm:wintner}
Let $f=\mathbf1*h$, that is $f(n)=\sum_{d\mid n}h(d)$, and assume that $\sum_{d\ge1}\frac{|h(d)|}d<\infty$. Then $f$ has a mean value, and
\[
m(f)=\sum_{d\ge1}\frac{h(d)}d.
\]
If moreover $h$ is multiplicative, $m(f)=\prod_p\bigl(1+\frac{h(p)}p+\frac{h(p^2)}{p^2}+\cdots\bigr)$.
\end{theorem}
\begin{proof}
By Lemma~\ref{lem:convsum} with $g=\mathbf1$,
\[
\sum_{n\le x}f(n)=x\sum_{d\le x}\frac{h(d)}d-\sum_{d\le x}h(d)\Bigl\{\frac xd\Bigr\}.
\]
Let $S=\sum_{d\ge1}\frac{h(d)}d$. The first term is $xS-x\sum_{d>x}\frac{h(d)}d$, and the tail sum tends to $0$. The second term is at most
$\sum_{d\le x}|h(d)|$ in modulus, and this is $o(x)$: for $1\le D\le x$, using $1\le\frac xd$ for $d\le x$,
\[
\begin{aligned}\frac1x\sum_{d\le x}|h(d)|&\le\frac1x\sum_{d\le D}|h(d)|+\sum_{D<d\le x}\frac{|h(d)|}d\\&\le\frac1x\sum_{d\le D}|h(d)|+\sum_{d>D}\frac{|h(d)|}d.\end{aligned}
\]
Let $x\to\infty$ with $D$ fixed: the $\limsup$ of the left side is at most $\sum_{d>D}\frac{|h(d)|}d$, which tends to $0$ as $D\to\infty$. Hence
$\frac1x\sum_{n\le x}f(n)\to S$. If $h$ is multiplicative, so is $\frac hd$, and $S$ is the stated Euler product by Volume~5, Theorem~12.3, applied at
$s=1$ to the absolutely convergent series $\sum\frac{h(d)}d$.
\end{proof}

The argument fails for $\tau=\mathbf1*\mathbf1$, where $\sum\frac1d$ diverges: the mean of $\tau$ is infinite, and the partial sums grow
like $x\log x$. To apply the theorem, we need the following identity for the squarefree indicator.

\begin{lemma}\label{lem:mu2}
For every $n\ge1$, $\mu^2(n)=\sum_{k^2\mid n}\mu(k)$.
\end{lemma}
\begin{proof}
Write $n=a^2b$ with $b$ squarefree; this is possible in exactly one way (put in $a^2$ the largest even power of each prime of $n$).
Then $k^2\mid n$ if and only if $k\mid a$: if $p^j\,\|\,k$ and $p^e\,\|\,n$, the condition $2j\le e$ means $j\le\lfloor e/2\rfloor$, the exponent of $p$
in $a$. So the sum is $\sum_{k\mid a}\mu(k)$, which is $1$ if $a=1$ and $0$ otherwise (Volume~2, Theorem~4.4). And $a=1$ exactly when
$n$ is squarefree.
\end{proof}

\begin{corollary}\label{cor:means}
\begin{enumerate}[label=(\alph*)]
\item The squarefree integers have density $m(\mu^2)=\sum_k\frac{\mu(k)}{k^2}=\prod_p\bigl(1-\frac1{p^2}\bigr)=\frac1{\zeta(2)}=\frac6{\pi^2}$.
\item $m\bigl(\frac{\varphi(n)}n\bigr)=\prod_p\bigl(1-\frac1{p^2}\bigr)=\frac6{\pi^2}$.
\item $m\bigl(\frac{\sigma(n)}n\bigr)=\sum_d\frac1{d^2}=\zeta(2)=\frac{\pi^2}6$.
\end{enumerate}
\end{corollary}
\begin{proof}
(a) By Lemma~\ref{lem:mu2}, $\mu^2=\mathbf1*h$ with $h(k^2)=\mu(k)$ and $h(d)=0$ when $d$ is not a square. Then
$\sum_d\frac{|h(d)|}d\le\sum_k\frac1{k^2}<\infty$, and $\sum_d\frac{h(d)}d=\sum_k\frac{\mu(k)}{k^2}=\frac1{\zeta(2)}$ (Volume~5, Corollary~12.4).
(b) $\varphi=\mu*\operatorname{id}$ gives $\frac{\varphi(n)}n=\sum_{d\mid n}\frac{\mu(d)}d$, so $h(d)=\frac{\mu(d)}d$ and $\sum_d\frac{h(d)}d=\sum_d\frac{\mu(d)}{d^2}$.
(c) $\frac{\sigma(n)}n=\sum_{d\mid n}\frac{n/d}n=\sum_{d\mid n}\frac1d$, so $h(d)=\frac1d$.
The Euler products come from Theorem~\ref{thm:wintner} or from the Euler product of $\zeta(2)$.
\end{proof}

\begin{beauty}
Read Corollary~\ref{cor:means}(b) probabilistically. Choose $n\le x$ at random. The events ``$p\mid n$'' for different primes $p$ have
probabilities close to $\frac1p$ and are nearly independent (Chapter~1, Exercise~\ref{ex:sieveprob}). The factor of $\frac{\varphi(n)}n=\prod_{p\mid n}(1-\frac1p)$
at $p$ is $1-\frac1p$ if $p\mid n$ and $1$ otherwise, with expectation $\frac1p(1-\frac1p)+(1-\frac1p)=1-\frac1{p^2}$. If expectation and product
commuted, as for independent variables (Volume~4), the mean would be $\prod_p(1-\frac1{p^2})$, which is right. Wintner's theorem says
that this naive model is exact whenever $h$ is small: only the small primes matter, and for them independence is true in the limit.
Compare with the factor $2e^{-\gamma}$ of Chapter~1, where the primes up to $\sqrt x$ all matter and the model fails. Chapter~9 makes
this model, due to Kubilius, into a tool.
\end{beauty}

\section{Error terms}

A mean value is a promise about the far future: it says where the averages end up, not how long they take to get there. Mean values say nothing about the speed of convergence. For explicit error terms we need three elementary estimates.

\begin{lemma}\label{lem:elemsums}
For real $x\ge1$:
\begin{enumerate}[label=(\alph*)]
\item $\sum_{n\le x}\frac1n=\log x+\gamma+O\bigl(\frac1x\bigr)$;
\item for $s>1$, $0\le\sum_{n>x}\frac1{n^s}\le\frac{2^{s-1}}{s-1}x^{1-s}$;
\item $\bigl|\sum_{n\le x}n-\frac{x^2}2\bigr|\le\frac x2$.
\end{enumerate}
\end{lemma}
\begin{proof}
(a) Let $N=\lfloor x\rfloor$. By Volume~5, Corollary~13.6, $\sum_{n\le N}\frac1n=\log N+\gamma+O(\frac1N)$, and $0\le\log x-\log N\le\frac{x-N}N\le\frac1N\le\frac2x$.
(b) The sum runs over $n\ge N+1$, and $n^{-s}\le\int_{n-1}^nt^{-s}\dd t$, so it is at most $\int_N^\infty t^{-s}\dd t=\frac{N^{1-s}}{s-1}\le\frac{(x/2)^{1-s}}{s-1}$, as $N\ge\frac x2$.
(c) $\sum_{n\le x}n=\frac{N(N+1)}2$, and $(x-1)x\le N(N+1)\le x(x+1)$.
\end{proof}

\begin{theorem}\label{thm:errorterms}
For $x\ge2$:
\begin{enumerate}[label=(\alph*)]
\item (Mertens, 1874) $\displaystyle\sum_{n\le x}\varphi(n)=\frac3{\pi^2}x^2+O(x\log x)$;
\item $\displaystyle\sum_{n\le x}\sigma(n)=\frac{\pi^2}{12}x^2+O(x\log x)$;
\item the number $Q(x)$ of squarefree integers $n\le x$ is $\frac6{\pi^2}x+O(\sqrt x)$.
\end{enumerate}
\end{theorem}
\begin{proof}
Let $S(y)=\sum_{m\le y}m=\frac{y^2}2+O(y)$ for $y\ge1$ (Lemma~\ref{lem:elemsums}(c)).

(a) By Lemma~\ref{lem:convsum} with $\varphi=\mu*\operatorname{id}$,
\[
\begin{aligned}\sum_{n\le x}\varphi(n)&=\sum_{d\le x}\mu(d)S\Bigl(\frac xd\Bigr)\eqbr=\frac{x^2}2\sum_{d\le x}\frac{\mu(d)}{d^2}+O\Bigl(x\sum_{d\le x}\frac1d\Bigr).\end{aligned}
\]
The error term is $O(x\log x)$ by Lemma~\ref{lem:elemsums}(a), and $\sum_{d\le x}\frac{\mu(d)}{d^2}=\frac6{\pi^2}+O(\frac1x)$ by Lemma~\ref{lem:elemsums}(b) with $s=2$.

(b) Here $\sigma=\mathbf1*\operatorname{id}$, so $\sum_{n\le x}\sigma(n)=\sum_{d\le x}S(\frac xd)=\frac{x^2}2\sum_{d\le x}\frac1{d^2}+O(x\log x)$, and $\sum_{d\le x}\frac1{d^2}=\frac{\pi^2}6+O(\frac1x)$.

(c) By Lemma~\ref{lem:mu2} and Lemma~\ref{lem:convsum} (the function $h$ of Corollary~\ref{cor:means}(a) vanishes off the squares),
\[
\begin{aligned}Q(x)&=\sum_{k\le\sqrt x}\mu(k)\Bigl\lfloor\frac x{k^2}\Bigr\rfloor\eqbr=x\sum_{k\le\sqrt x}\frac{\mu(k)}{k^2}+O(\sqrt x).\end{aligned}
\]
By Lemma~\ref{lem:elemsums}(b), $\sum_{k\le\sqrt x}\frac{\mu(k)}{k^2}=\frac6{\pi^2}+O(\frac1{\sqrt x})$, and the result follows.
\end{proof}

Counting coprime pairs reduces to $\varphi$. Figure~\ref{fig:coprime} shows them as the lattice points visible from the origin.

\begin{corollary}\label{cor:coprimepairs}
The number $C(x)$ of pairs $(a,b)\in\{1,\dots,\lfloor x\rfloor\}^2$ with $\gcd(a,b)=1$ satisfies $C(x)=\frac6{\pi^2}x^2+O(x\log x)$ for $x\ge2$.
\end{corollary}
\begin{proof}
For $b\ge2$, the pairs $(a,b)$ with $a<b$ and $\gcd(a,b)=1$ number $\varphi(b)$, and the pair $(b,b)$ is not coprime. For $b=1$ only $(1,1)$
qualifies, and $\varphi(1)=1$. By the symmetry $(a,b)\mapsto(b,a)$, $C(x)=2\sum_{b\le x}\varphi(b)-1$. Apply Theorem~\ref{thm:errorterms}(a).
\end{proof}

\begin{figure}[ht]
\centering
\fig{ch02-coprime}
\caption{The pairs $(a,b)$ with $1\le a,b\le24$: coprime pairs filled, the others hollow. A point is coprime exactly when it is the first
lattice point on its ray from the origin, so the filled points are those visible from $(0,0)$; on the coral ray only $(3,2)$ is. Here
$359$ of the $576$ points are visible, a proportion $0.623$; the limit is $\frac6{\pi^2}=0.6079\ldots$}
\label{fig:coprime}
\end{figure}

\begin{remark}
These error terms are far from the truth, but improving them is hard. For $\varphi$, the best known error is
$O\bigl(x(\log x)^{2/3}(\log\log x)^{4/3}\bigr)$ (Walfisz, 1963). For $Q(x)$, no bound $O(x^{\theta})$ with $\theta<\frac12$ is known without hypothesis;
under the Riemann hypothesis, Montgomery and Vaughan (1981) obtained $O(x^{9/28+\varepsilon})$. The reason is the same in both
cases: the error is governed by the sums $\sum_{n\le y}\mu(n)$, which are controlled by the zeros of $\zeta$ (Chapters~5 and~7).
These results are quoted for culture; see the sources at the end of the chapter.
\end{remark}

\section{The hyperbola method}

Dirichlet's best trick is a matter of looking at a picture. Volume~2, Theorem~4.7 counted the lattice points under the hyperbola $dm=x$ by symmetry. The same idea works for any
convolution.

\begin{lemma}[Dirichlet's hyperbola method]\label{lem:hyperbola}
Let $f=g*h$, with $G(y)=\sum_{m\le y}g(m)$ and $H(y)=\sum_{m\le y}h(m)$. For $1\le y\le x$,
\[
\begin{aligned}\sum_{n\le x}f(n)&=\sum_{d\le y}g(d)H\Bigl(\frac xd\Bigr)+\sum_{m\le x/y}h(m)G\Bigl(\frac xm\Bigr)\eqbr-G(y)H\Bigl(\frac xy\Bigr).\end{aligned}
\]
\end{lemma}
\begin{proof}
The left side is $\sum g(d)h(m)$ over the pairs with $dm\le x$. Each such pair has $d\le y$ or $m\le\frac xy$, since otherwise $dm>x$. The
first sum covers the pairs with $d\le y$, the second those with $m\le\frac xy$, and the pairs with both, counted twice, contribute
$\sum_{d\le y}g(d)\sum_{m\le x/y}h(m)=G(y)H(\frac xy)$.
\end{proof}

For $g=h=\mathbf1$ and $y=\sqrt x$ this is Volume~2's formula $\sum_{n\le x}\tau(n)=2\sum_{d\le\sqrt x}\lfloor\frac xd\rfloor-\lfloor\sqrt x\rfloor^2$, which gave
\[
\begin{aligned}D(x)&=\sum_{n\le x}\tau(n)\eqbr=x\log x+(2\gamma-1)x+\Delta(x),\end{aligned}
\]
with $|\Delta(x)|\le C\sqrt x$ for $x\ge1$ and an absolute constant $C$. Iterating gives the three-fold divisor function $\tau_3=\mathbf1*\mathbf1*\mathbf1=\tau*\mathbf1$, the number of
ways to write $n=abc$ as an ordered product of three positive integers. We need one more elementary sum.

\begin{lemma}\label{lem:logsum}
$\displaystyle\sum_{m\le x}\frac{\log m}m=\frac12\log^2x+O(1)$ for $x\ge1$.
\end{lemma}
\begin{proof}
Apply Lemma~\ref{lem:abelint} with $a_m=1$, $A(t)=\lfloor t\rfloor$, $f(t)=\frac{\log t}t$ and $y=1$, and note $f(1)=0$. Writing $\lfloor t\rfloor=t-\{t\}$
and integrating $tf'(t)$ by parts,
\[
\begin{aligned}\sum_{m\le x}\frac{\log m}m&=\lfloor x\rfloor f(x)-xf(x)+\int_1^xf(t)\dd t\eqbr+\int_1^x\{t\}f'(t)\dd t.\end{aligned}
\]
Here $\int_1^xf=\frac12\log^2x$, $|\lfloor x\rfloor f(x)-xf(x)|\le f(x)\le\frac1e$, and $|f'(t)|=\frac{|1-\log t|}{t^2}$ is integrable on $[1,\infty[$.
\end{proof}

\begin{theorem}\label{thm:tau3}
For $x\ge2$,
\[
\sum_{n\le x}\tau_3(n)=\frac12x\log^2x+(3\gamma-1)x\log x+O(x).
\]
\end{theorem}
\begin{proof}
By Lemma~\ref{lem:convsum} with $\tau_3=\mathbf1*\tau$, $\sum_{n\le x}\tau_3(n)=\sum_{m\le x}D(\frac xm)$, and $\frac xm\ge1$. Insert the formula for $D$:
\[
\begin{aligned}\sum_{m\le x}D\Bigl(\frac xm\Bigr)&=x\sum_{m\le x}\frac{\log x-\log m}m\eqbr+(2\gamma-1)x\sum_{m\le x}\frac1m+\sum_{m\le x}\Delta\Bigl(\frac xm\Bigr).\end{aligned}
\]
By Lemma~\ref{lem:elemsums}(a) and Lemma~\ref{lem:logsum}, the first term is $x\log x(\log x+\gamma)-\frac12x\log^2x+O(x)$, the second is
$(2\gamma-1)x\log x+O(x)$, and the third is at most $C\sqrt x\sum_{m\le x}m^{-1/2}\le C\sqrt x\cdot2\sqrt x$, since $\sum_{m\le x}m^{-1/2}\le\int_0^xt^{-1/2}\dd t$. Adding up gives
$\frac12x\log^2x+(3\gamma-1)x\log x+O(x)$.
\end{proof}

The pattern continues:
\[
\sum_{n\le x}\tau_k(n)=xP_k(\log x)+O(x^{1-1/k}\log^{k-2}x),
\]
where $P_k$ is a polynomial of degree $k-1$ with leading
coefficient $\frac1{(k-1)!}$ (Exercise~\ref{ex:tauk} proves the main term). Numerically, one finds the constant term $0.4863\ldots$ of
$P_3$; it involves the Stieltjes constant $\gamma_1$, the next coefficient after $\gamma$ in the expansion of $\zeta$ at $1$.

\section{The divisor problem}

Now we reach the edge of the map. How large is $\Delta(x)$ really? Figure~\ref{fig:divisorerr} plots $\Delta(x)/x^{1/4}$ and shows a bounded-looking oscillation, far below Dirichlet's
$O(\sqrt x)$.

\begin{figure}[ht]
\centering
\fig{ch02-divisorerr}
\caption{$\Delta(x)/x^{1/4}$, where $\Delta(x)=\sum_{n\le x}\tau(n)-x\log x-(2\gamma-1)x$, plotted against $\sqrt x$ for $400\le x\le1600$. It oscillates
with a root mean square close to $1$, as Cram\'er's mean-square theorem predicts; Dirichlet's bound only says that the curve stays below a multiple of
$\sqrt x/x^{1/4}=x^{1/4}$, which is between $4.5$ and $6.3$ here.}
\label{fig:divisorerr}
\end{figure}

\begin{remark}[for culture]
Finding the true order of $\Delta(x)$ is \emph{Dirichlet's divisor problem}, open since 1849. Vorono\"\i{} (1903) proved $\Delta(x)=O(x^{1/3}\log x)$
with a much finer analysis of the lattice points near the hyperbola; Huxley (2003) reached the exponent $\frac{131}{416}=0.3149\ldots$. In
the other direction, Hardy (1916) proved that $\Delta(x)$ is not $O(x^{1/4})$, and Cram\'er (1922) that
\[
\begin{gathered}\int_1^X\Delta(x)^2\dd x\sim cX^{3/2},\qq c=\frac1{6\pi^2}\sum_{n\ge1}\frac{\tau(n)^2}{n^{3/2}}=0.6543\ldots\end{gathered}
\]
so that $\Delta(x)$ is of size $x^{1/4}$ on average. The conjecture is $\Delta(x)=O(x^{1/4+\varepsilon})$ for every $\varepsilon>0$. The same story, with essentially the
same exponents, holds for Gauss's circle problem (Exercise~\ref{ex:circle}); there Sierpi\'nski proved $O(x^{1/3})$ in 1906. Proofs are in Titchmarsh's book, Chapter~XII.
\end{remark}

The average of $\tau$ hides another surprise, which the think box was after. The mean of $\tau(n)$ for $n\le10^6$ is about $14$, but
the median is $8$, and only $34\%$ of these integers have $14$ divisors or more. The average is pulled up by rare integers with very
many divisors. Chapter~9 shows that for almost all $n$, $\tau(n)=(\log n)^{\log2+o(1)}$, and $(\log10^6)^{\log2}\approx6.2$: the typical
integer has far fewer divisors than the average one. Still, no integer has very many divisors, compared with its size
(Exercise~\ref{ex:taubound}).

\section{Exercises}

\begin{exercise}[$\star$]\label{ex:sigmaovern}
Show that $\sum_{n\le x}\frac{\sigma(n)}n=\frac{\pi^2}6x+O(\log x)$ and $\sum_{n\le x}\frac{\varphi(n)}n=\frac6{\pi^2}x+O(\log x)$ for $x\ge2$.
\end{exercise}

\begin{exercise}[$\star\star$]\label{ex:nphi}
Show that $\frac n{\varphi(n)}=\sum_{d\mid n}\frac{\mu^2(d)}{\varphi(d)}$, and deduce that $\frac n{\varphi(n)}$ has mean value
$\prod_p\bigl(1+\frac1{p(p-1)}\bigr)=\frac{\zeta(2)\zeta(3)}{\zeta(6)}=1.9435\ldots$
\end{exercise}

\begin{exercise}[$\star\star$]\label{ex:coprimek}
\emph{Signature exercise.} Let $k\ge2$ and let $C_k(x)$ be the number of $k$-tuples $(a_1,\dots,a_k)\in\{1,\dots,\lfloor x\rfloor\}^k$ with
$\gcd(a_1,\dots,a_k)=1$. Show that $C_k(x)=\sum_{d\le x}\mu(d)\lfloor\frac xd\rfloor^k$, and that
\[
\frac{C_k(x)}{x^k}=\frac1{\zeta(k)}+\begin{cases}O\bigl(\frac{\log x}x\bigr)&k=2,\\O\bigl(\frac1x\bigr)&k\ge3.\end{cases}
\]
So three random integers are coprime with probability $\frac1{\zeta(3)}=0.8319\ldots$
\end{exercise}

\begin{exercise}[$\star\star$]\label{ex:circle}
\emph{Gauss's circle problem.} Let $N(r)$ be the number of $(a,b)\in\Z^2$ with $a^2+b^2\le r^2$. By placing a unit square centred at each
of these points, show that $|N(r)-\pi r^2|\le\pi(\sqrt2\,r+\frac12)$ for $r\ge1$. Deduce that $\sum_{0\le n\le x}r_2(n)=\pi x+O(\sqrt x)$, where $r_2(n)$
is the number of $(a,b)\in\Z^2$ with $a^2+b^2=n$, so that $r_2$ has mean value $\pi$.
\end{exercise}

\begin{exercise}[$\star\star$]\label{ex:twoomega}
Show that $\sum_{n\le x}2^{\omega(n)}=\frac6{\pi^2}x\log x+O(x)$ for $x\ge2$, where $\omega(n)$ is the number of distinct prime factors of $n$.
\end{exercise}

\begin{exercise}[$\star\star$]\label{ex:taubound}
\begin{enumerate}[label=(\alph*)]
\item Show that $\tau(n)\le\sqrt{3n}$ for every $n\ge1$, with equality only for $n=12$.
\item Show that for every $\varepsilon\in\left]0,1\right]$ and every $n\ge1$, $\tau(n)\le(\varepsilon\log2)^{-2^{1/\varepsilon}}n^\varepsilon$.
\end{enumerate}
\end{exercise}

\begin{exercise}[$\star\star$]\label{ex:tauk}
Let $\tau_k=\mathbf1^{*k}$ be the number of ordered factorisations $n=a_1\cdots a_k$. Show by induction on $k\ge1$ that
$\sum_{n\le x}\tau_k(n)=\frac1{(k-1)!}x\log^{k-1}x+O_k(x\log^{k-2}x)$ for $x\ge3$, where for $k=1$ the error term is read as $O(1)$.
\end{exercise}

\begin{exercise}[$\star\star$]\label{pb:phiovern}
\emph{The law of $\varphi(n)/n$.}
\begin{enumerate}[label=(\alph*)]
\item Show that for every integer $k\ge1$,
\[
\frac1N\sum_{n\le N}\Bigl(\frac{\varphi(n)}n\Bigr)^k\longrightarrow\prod_p\Bigl(1-\frac1p+\frac1p\Bigl(1-\frac1p\Bigr)^k\Bigr).
\]
\item Deduce that $\frac1N\sum_{n\le N}g\bigl(\frac{\varphi(n)}n\bigr)$ has a limit for every continuous $g:[0,1]\to\R$: the values
$\varphi(n)/n$ have a limiting distribution.
\item Show that at least half of all integers, asymptotically, have $\varphi(n)\le\frac n2$, although the mean of
$\varphi(n)/n$ is $\frac6{\pi^2}>\frac12$.
\end{enumerate}
\end{exercise}

\begin{exercise}[$\star\star$]\label{pb:visible}
\emph{Visible points.} A point $(a,b)\in\Z^2\setminus\{0\}$ is \emph{visible} from the origin if no other lattice point lies on the
segment joining them.
\begin{enumerate}[label=(\alph*)]
\item Show that $(a,b)$ is visible if and only if $\gcd(a,b)=1$.
\item Let $V(R)$ be the number of visible points in the disc of radius $R$. Show that $V(R)=\frac6\pi R^2+O(R\log R)$.
\item Show that the visible points are equidistributed in angle: the proportion of them in a sector of angle $\theta$ tends to
$\frac\theta{2\pi}$.
\end{enumerate}
\end{exercise}

\begin{exercise}[$\star\star$]\label{pb:farey}
\emph{Farey fractions.} The Farey sequence $\mathcal F_N$ lists the fractions $\frac hk$ in $[0,1]$ in lowest terms with $1\le k\le N$,
in increasing order.
\begin{enumerate}[label=(\alph*)]
\item Show that $|\mathcal F_N|=1+\sum_{k\le N}\varphi(k)=\frac3{\pi^2}N^2+O(N\log N)$.
\item Show that two consecutive fractions $\frac hk<\frac{h'}{k'}$ of $\mathcal F_N$ satisfy $h'k-hk'=1$ and $k+k'>N$.
\item Deduce that $\sum\frac1{kk'}=1$, the sum running over consecutive pairs, and that the gaps of $\mathcal F_N$ all lie between
$\frac1{N(N-1)}$ and $\frac1N$.
\end{enumerate}
\end{exercise}

\begin{exercise}[$\star\star$]\label{pb:orderedfact}
\emph{Ordered factorisations (Kalmár, 1931).} Let $H(1)=1$ and, for $n\ge2$, let $H(n)$ be the number of ways of writing $n$ as an
ordered product of factors $>1$ (so $H(12)=8$). Let $\rho>1$ be the solution of $\zeta(\rho)=2$; numerically $\rho=1.7286\ldots$
\begin{enumerate}[label=(\alph*)]
\item Show that $\sum_nH(n)n^{-s}=\frac1{2-\zeta(s)}$ for real $s>\rho$.
\item Compute $H(p^a)$ for a prime $p$, and $H(p_1\cdots p_k)$ for distinct primes, for $k\le4$.
\item Show that $H(n)\le n^\rho$ for every $n\ge1$.
\item Show that $\sum_{n\le x}H(n)\le\frac{x^\sigma}{2-\zeta(\sigma)}$ for every $\sigma>\rho$ and $x\ge1$.
\end{enumerate}
\end{exercise}

\begin{exercise}[$\star\star$]\label{pb:fracmean}
\emph{Remainders on average.}
\begin{enumerate}[label=(\alph*)]
\item Show that $\displaystyle\sum_{n\le x}\Bigl\{\frac xn\Bigr\}=(1-\gamma)x+O(\sqrt x)$ for $x\ge1$.
\item For an integer $N\ge2$, let $N\bmod n$ be the remainder of the division of $N$ by $n$. Show that
\[
\sum_{n=1}^N(N\bmod n)=\Bigl(1-\frac{\pi^2}{12}\Bigr)N^2+O(N\log N).
\]
\item Explain both constants as integrals: $1-\gamma=\int_0^1\{\frac1t\}\dd t$ and $2-\frac{\pi^2}6=\int_0^12t\{\frac1t\}\dd t$.
\end{enumerate}
\end{exercise}

\begin{exercise}[$\star\star$]\label{pb:radical}
\emph{The average radical.} The radical $\operatorname{rad}(n)$ is the product of the distinct primes dividing $n$, so $\operatorname{rad}(72)=6$. It is the quantity in
the $abc$ conjecture of Masser and Oesterl\'e. Let $C=\prod_p\bigl(1-\frac1{p(p+1)}\bigr)=0.70444\ldots$
\begin{enumerate}[label=(\alph*)]
\item Write $\frac{\operatorname{rad}(n)}n=\sum_{d\mid n}g(d)$ and compute $g$. Show that $g(d)=0$ unless $d$ is \emph{squarefull} (every prime factor of $d$ appears at
least squared), and that $|g(d)|\le1$.
\item Show that the number of squarefull $d\le y$ is at most $\zeta(\frac32)\sqrt y$.
\item Deduce that $\sum_{n\le x}\operatorname{rad}(n)=\frac C2x^2+O(x^{3/2})$, and that the mean value of $\frac{\operatorname{rad}(n)}n$ is $C$.
\end{enumerate}
\end{exercise}

\section{Hints}

\begin{hint}{ex:sigmaovern}
$\frac{\sigma(n)}n=\sum_{d\mid n}\frac1d$ and $\frac{\varphi(n)}n=\sum_{d\mid n}\frac{\mu(d)}d$; swap the sums.
\end{hint}
\begin{hint}{ex:nphi}
Both sides are multiplicative; compare them at prime powers. Then apply Theorem~\ref{thm:wintner}.
\end{hint}
\begin{hint}{ex:coprimek}
$\mathbf 1_{\gcd=1}=\sum_{d\mid\gcd}\mu(d)$, and $y^k-\lfloor y\rfloor^k\le ky^{k-1}$ for $y\ge1$.
\end{hint}
\begin{hint}{ex:circle}
The square centred at $(a,b)$ lies within distance $\frac{\sqrt2}2$ of $(a,b)$. Every point of the plane lies in the square of its
nearest lattice point.
\end{hint}
\begin{hint}{ex:twoomega}
$2^{\omega}=\mu^2*\mathbf1$ (Volume~2, Exercise~4.2). Then use Theorem~\ref{thm:errorterms}(c) and Abel summation.
\end{hint}
\begin{hint}{ex:taubound}
Write $\frac{\tau(n)}{n^\varepsilon}=\prod_{p^k\|n}\frac{k+1}{p^{k\varepsilon}}$ and bound each factor. For which $p$ can a factor exceed $1$?
\end{hint}
\begin{hint}{ex:tauk}
$\tau_{k+1}=\tau_k*\mathbf1$; use Lemma~\ref{lem:convsum} with $h=\mathbf1$ and the induction hypothesis at $\frac xm$, and compare
$\sum_{m\le x}\frac{\log^{j}(x/m)}m$ with $\int_1^x\frac{\log^j(x/t)}t\dd t=\frac{\log^{j+1}x}{j+1}$.
\end{hint}

\begin{hint}{pb:phiovern}
(a) $(\varphi(n)/n)^k$ is multiplicative and close to $1$ at large primes: use Wintner's theorem (Volume~6, Theorem~2.2).
(b) Weierstrass.
\end{hint}

\begin{hint}{pb:visible}
(b) Every non-zero lattice point is a multiple $d\cdot v$ of a unique visible point $v$; invert with Möbius.
\end{hint}

\begin{hint}{pb:farey}
(b) Solve $kx-hy=1$ with $N-k<y\le N$, and show that $\frac xy$ is the successor of $\frac hk$.
\end{hint}

\begin{hint}{pb:orderedfact}
(a) $H=\sum_k(\zeta-1)^k$. (c) Strong induction, with $H(n)=\sum_{d\mid n,\,d<n}H(d)$. (d) Rankin's trick.
\end{hint}

\begin{hint}{pb:fracmean}
(a) $\sum_{n\le x}\lfloor\frac xn\rfloor=\sum_{m\le x}\tau(m)$. (b) $\sum_{n\le N}n\lfloor\frac Nn\rfloor=\sum_{m\le N}\sigma(m)$.
\end{hint}

\begin{hint}{pb:radical}
(a) $g=f*\mu$ with $f(p^k)=p^{1-k}$; compute $g(p^k)$. (b) Every squarefull number is $a^2b^3$. (c) $\sum_{n\le x}\operatorname{rad}(n)=\sum_dg(d)\,d\sum_{m\le x/d}m$.
\end{hint}

\section{Solutions}

\begin{solution}{ex:sigmaovern}
\idea Each function is $\mathbf1*h$ with $h(d)=\frac1d$ or $\frac{\mu(d)}d$. Lemma~\ref{lem:convsum} gives the main term $x\sum\frac{h(d)}d$ and an error
$\sum_{d\le x}|h(d)|$, which here is only $O(\log x)$.
\writeup By Lemma~\ref{lem:convsum}, $\sum_{n\le x}\frac{\sigma(n)}n=\sum_{d\le x}\frac1d\lfloor\frac xd\rfloor=x\sum_{d\le x}\frac1{d^2}-\sum_{d\le x}\frac1d\{\frac xd\}$. By Lemma~\ref{lem:elemsums},
the first term is $\frac{\pi^2}6x+O(1)$ and the second is between $0$ and $\log x+1$. The same computation with $\frac{\mu(d)}d$ gives
$x\sum_{d\le x}\frac{\mu(d)}{d^2}+O(\log x)=\frac6{\pi^2}x+O(\log x)$. These are Corollary~\ref{cor:means}(b) and~(c) with error terms.
\end{solution}

\begin{solution}{ex:nphi}
\idea $\frac n{\varphi(n)}=\prod_{p\mid n}\frac p{p-1}$ depends only on the primes dividing $n$, which is exactly what a sum over squarefree divisors
produces. The function $h=\frac{\mu^2}\varphi$ is small enough for Wintner.
\writeup Both sides are multiplicative functions of $n$ (the right side is $\mathbf1*h$ with $h=\frac{\mu^2}\varphi$ multiplicative; Volume~2, Proposition~4.3). At
$n=p^k$ with $k\ge1$, the left side is $\frac{p^k}{p^{k-1}(p-1)}=\frac p{p-1}$, and the right side is $1+\frac1{p-1}=\frac p{p-1}$, since $\mu^2(p^j)=0$ for $j\ge2$.
So they agree. Next, $\frac{h(d)}d\ge0$ is multiplicative, with $\frac{h(p)}p=\frac1{p(p-1)}$ and $h(p^j)=0$ for $j\ge2$. Every $d\le x$ with $h(d)\ne0$
is a squarefree product of primes $\le x$, so $\sum_{d\le x}\frac{h(d)}d\le\prod_{p\le x}\bigl(1+\frac1{p(p-1)}\bigr)\le\exp\sum_p\frac1{p(p-1)}<\infty$. Theorem~\ref{thm:wintner}
applies and gives the mean value $\prod_p(1+\frac1{p(p-1)})$. Finally, the Euler factor of $\frac{\zeta(2)\zeta(3)}{\zeta(6)}$ at $p$ is
\[
\begin{aligned}\frac{1-p^{-6}}{(1-p^{-2})(1-p^{-3})}&=\frac{1+p^{-3}}{1-p^{-2}}\eqbr=\frac{(p+1)(p^2-p+1)}{p(p-1)(p+1)}\eqbr=1+\frac1{p(p-1)}.\end{aligned}
\]
Numerically $1.94359\ldots$; one finds numerically $1.94358\ldots$ for the average up to $10^6$. So $n/\varphi(n)$ is about $2$ on average, while
Corollary~\ref{cor:means}(b) gave $0.608$ for $\varphi(n)/n$: the mean of a reciprocal is not the reciprocal of the mean.
\end{solution}

\begin{solution}{ex:coprimek}
\idea As in Volume~2, Exercise~4.8, Möbius inversion turns the coprimality condition into a sum over common divisors $d$, and the $k$-tuples
of multiples of $d$ are easy to count. The new point is the error term, which is smaller for $k\ge3$ because $\sum d^{1-k}$ converges.
\writeup Since $\sum_{d\mid g}\mu(d)$ is $1$ if $g=1$ and $0$ otherwise,
\[
C_k(x)=\sum_{a_1,\dots,a_k}\sum_{d\mid a_1,\dots,d\mid a_k}\mu(d)=\sum_{d\le x}\mu(d)\lfloor\frac xd\rfloor^k,
\]
because
each $a_i$ ranges over the $\lfloor\frac xd\rfloor$ multiples of $d$ up to $x$. For $y\ge1$, the mean value theorem gives $0\le y^k-\lfloor y\rfloor^k\le ky^{k-1}(y-\lfloor y\rfloor)\le ky^{k-1}$.
So
\[
\begin{aligned}C_k(x)&=x^k\sum_{d\le x}\frac{\mu(d)}{d^k}+O\Bigl(kx^{k-1}\sum_{d\le x}\frac1{d^{k-1}}\Bigr).\end{aligned}
\]
By Lemma~\ref{lem:elemsums}(b) and Volume~5, Corollary~12.4, $\sum_{d\le x}\frac{\mu(d)}{d^k}=\frac1{\zeta(k)}+O(x^{1-k})$, which contributes $O(x)$. The sum
$\sum_{d\le x}d^{1-k}$ is $O(\log x)$ for $k=2$ and bounded by $\zeta(2)$ for $k\ge3$. Dividing by $x^k$ gives the claim. For $x=1000$ one finds numerically
$0.60838$, $0.83205$, $0.92399$ for $k=2,3,4$, against $0.60793$, $0.83191$, $0.92394$.
\end{solution}

\begin{solution}{ex:circle}
\idea Area counts lattice points up to a boundary layer: each point owns a unit square, and the squares of the points in the disc fill the
disc up to a strip of width $\frac{\sqrt2}2$, the half-diagonal of a square, along the circle.
\writeup Let $K(a,b)=[a-\frac12,a+\frac12]\times[b-\frac12,b+\frac12]$; these squares have area $1$ and disjoint interiors, and every point of $K(a,b)$ is within
$\frac{\sqrt2}2$ of $(a,b)$. If $a^2+b^2\le r^2$, then $K(a,b)$ lies in the disc of radius $r+\frac{\sqrt2}2$, so $N(r)\le\pi(r+\frac{\sqrt2}2)^2=\pi r^2+\pi(\sqrt2r+\frac12)$.
Conversely, let $|z|\le r-\frac{\sqrt2}2$ and let $(a,b)$ be a lattice point nearest to $z$; then $z\in K(a,b)$ and $|(a,b)|\le|z|+\frac{\sqrt2}2\le r$. So the $N(r)$
squares cover the disc of radius $r-\frac{\sqrt2}2$, and $N(r)\ge\pi(r-\frac{\sqrt2}2)^2\ge\pi r^2-\pi\sqrt2r$. For the sum, $\sum_{0\le n\le x}r_2(n)=N(\sqrt x)$ counts
the same points. For $r=1000$ one finds numerically $N(r)-\pi r^2=-43.65$, much less than the bound $4444$: the true error is believed to be
$O(r^{1/2+\varepsilon})$.
\end{solution}

\begin{solution}{ex:twoomega}
\idea $2^{\omega(n)}$ counts the squarefree divisors of $n$, so its sum counts pairs (squarefree $d$, multiple of $d$). The squarefree $d$ have
density $\frac6{\pi^2}$, and summing $\frac xd$ over them gives $\frac6{\pi^2}x\log x$.
\writeup By Volume~2, Exercise~4.2, $2^{\omega(n)}=\sum_{d\mid n}\mu^2(d)$. Lemma~\ref{lem:convsum} gives
\[
\sum_{n\le x}2^{\omega(n)}=\sum_{d\le x}\mu^2(d)\Bigl\lfloor\frac xd\Bigr\rfloor=x\sum_{d\le x}\frac{\mu^2(d)}d+O(x).
\]
Abel summation (Lemma~\ref{lem:abelint} with $a_d=\mu^2(d)$, $A=Q$, $f(t)=\frac1t$) gives $\sum_{d\le x}\frac{\mu^2(d)}d=\frac{Q(x)}x+\int_1^x\frac{Q(t)}{t^2}\dd t$. By
Theorem~\ref{thm:errorterms}(c), which also holds for $1\le t<2$ since $Q(t)=1$ there, $Q(t)=\frac6{\pi^2}t+O(\sqrt t)$, so the integral is
$\frac6{\pi^2}\log x+O\bigl(\int_1^\infty t^{-3/2}\dd t\bigr)=\frac6{\pi^2}\log x+O(1)$, and $\frac{Q(x)}x\le1$. Multiplying by $x$ gives the result. The constant
in the $O(x)$ can be computed: a direct computation verifies that the next term is $Cx$ with $C=\frac6{\pi^2}\bigl(2\gamma-1-2\frac{\zeta'(2)}{\zeta(2)}\bigr)=0.7869\ldots$
\end{solution}

\begin{solution}{ex:taubound}
\idea $\tau$ and $n^\varepsilon$ are both multiplicative, so compare them prime by prime. Large primes contribute factors at most $1$; only the
finitely many small primes can make $\frac{\tau(n)}{n^\varepsilon}$ big, and each of them only boundedly.
\writeup (a) $\frac{\tau(n)}{\sqrt n}=\prod_{p^k\|n}\frac{k+1}{p^{k/2}}$. For $p\ge5$ and $k\ge1$, $p^{k/2}\ge5^{k/2}>k+1$ (true for $k=1$, and the left side grows faster), so
these factors are $<1$. For $p=3$ the factors are $\frac2{\sqrt3}=1.155$, $\frac33=1$, then less than $1$; for $p=2$ they are $\frac2{\sqrt2}=1.414$,
$\frac32=1.5$, $\frac4{2\sqrt2}=1.414$, then decreasing. So $\frac{\tau(n)}{\sqrt n}\le\frac32\cdot\frac2{\sqrt3}=\sqrt3$, with equality only if $2^2\|n$, $3\|n$ and no other prime
divides $n$, that is $n=12$, where $\tau(12)=6=\sqrt{36}$.

(b) $\frac{\tau(n)}{n^\varepsilon}=\prod_{p^k\|n}\frac{k+1}{p^{k\varepsilon}}$. If $p\ge2^{1/\varepsilon}$, then $p^{k\varepsilon}\ge2^k\ge k+1$, so the factor is at most $1$. If $p<2^{1/\varepsilon}$, then
$p^{k\varepsilon}\ge2^{k\varepsilon}=e^{k\varepsilon\log2}\ge1+k\varepsilon\log2\ge(k+1)\varepsilon\log2$, using $\varepsilon\log2\le1$; so the factor is at most $\frac1{\varepsilon\log2}$. There are fewer
than $2^{1/\varepsilon}$ such primes, and $\frac1{\varepsilon\log2}>1$, so $\frac{\tau(n)}{n^\varepsilon}\le(\varepsilon\log2)^{-2^{1/\varepsilon}}$. The constant is enormous for small $\varepsilon$, but
the point is that $\tau(n)=O_\varepsilon(n^\varepsilon)$ for every $\varepsilon>0$. Optimising $\varepsilon$ as a function of $n$ gives Wigert's theorem
$\tau(n)\le2^{(1+o(1))\log n/\log\log n}$, attained by products of the first primes such as $720720=2^4\cdot3^2\cdot5\cdot7\cdot11\cdot13$.
\end{solution}

\begin{solution}{ex:tauk}
\idea Each convolution with $\mathbf1$ is a sum $\sum_{m\le x}T(\frac xm)$, and summing $\frac xm\log^{j}\frac xm$ over $m$ behaves like the integral
$x\int_1^x\log^j(\frac xt)\frac{\mathrm dt}t=\frac x{j+1}\log^{j+1}x$, which raises the power of the logarithm by one and divides by $j+1$.
\writeup Let $T_k(y)=\sum_{n\le y}\tau_k(n)$, $j=k-1$ and $L=\log x$. We prove the slightly more precise statement that for all $y\ge1$,
$T_k(y)=\frac{y\log^jy}{j!}+E_k(y)$ with $|E_k(y)|\le C_ky(1+\log y)^{j-1}$ if $k\ge2$, and $|E_1(y)|\le1$; for $y\ge3$ this is the claim, since
$1+\log y\le2\log y$. For $k=1$, $T_1(y)=\lfloor y\rfloor$. Suppose the statement for $k$. By Lemma~\ref{lem:convsum}, $T_{k+1}(x)=\sum_{m\le x}T_k(\frac xm)$, and
the errors add up to at most $\sum_{m\le x}C_k\frac xm(1+L)^{j-1}\le C_kx(1+L)^j$ (for $k=1$, at most $x$), which is of the required form. For the main term, let
$\phi(t)=\frac1t\log^j\frac xt$ on $[1,x]$. It is decreasing for $j=0$, and for $j\ge1$ it is decreasing too, since
$\phi'(t)=-\frac1{t^2}\log^{j-1}\frac xt\bigl(\log\frac xt+j\bigr)\le0$. For a decreasing function, $\int_1^{\lfloor x\rfloor+1}\phi\le\sum_{m\le x}\phi(m)\le\phi(1)+\int_1^x\phi$, with $\phi$
extended by $0$ beyond $x$; so the sum is $\int_1^x\phi+O(\phi(1))=\frac{L^{j+1}}{j+1}+O(L^j)$. Multiplying by $\frac x{j!}$ gives
$T_{k+1}(x)=\frac{xL^k}{k!}+O(xL^{k-1})$, the claim for $k+1$. For $k=3$ this is the main term of Theorem~\ref{thm:tau3}.
\end{solution}

\begin{solution}{pb:phiovern}
\idea $(\varphi(n)/n)^k$ is a multiplicative function equal to $\sum_{d\mid n}h(d)$ for a function $h$ that vanishes off squarefree
numbers and is $O(k/p)$ at primes. Wintner's theorem then gives the mean, and polynomials approximate every continuous function.

\writeup (a) Let $f_k(n)=(\varphi(n)/n)^k=\prod_{p\mid n}(1-\frac1p)^k$, multiplicative. Define $h$ multiplicative by
$h(p)=(1-\frac1p)^k-1$ and $h(p^j)=0$ for $j\ge2$. Then $\sum_{d\mid n}h(d)=\prod_{p\mid n}(1+h(p))=f_k(n)$. Since $|h(p)|\le\frac kp$,
\[
\sum_d\frac{|h(d)|}d\le\prod_p\Bigl(1+\frac k{p^2}\Bigr)<\infty.
\]
By Wintner's theorem (Volume~6, Theorem~2.2), $f_k$ has mean value $\sum_d\frac{h(d)}d=\prod_p\bigl(1+\frac{h(p)}p\bigr)$, which is the
stated product.

(b) By (a) and linearity, $\frac1N\sum_{n\le N}g(\varphi(n)/n)$ converges for every polynomial $g$. For $g$ continuous and
$\varepsilon>0$, Weierstrass's theorem (Volume~4, Theorem~4.5, via Bernstein polynomials) gives a polynomial $q$ with
$|g-q|\le\varepsilon$ on $[0,1]$. The averages of $g$ and of $q$ differ by at most $\varepsilon$, and those of $q$ form a Cauchy
sequence; so those of $g$ do too, and converge.

(c) Every even $n$ has $\varphi(n)/n\le\frac12$, and the even numbers have density $\frac12$. For $k=1$ the product in (a) is
$\prod_p(1-\frac1{p^2})=\frac6{\pi^2}$. The mean exceeds $\frac12$ because the odd numbers mostly have $\varphi(n)/n$ close to $1$.

\tools Volume~6, Chapter~2 (Wintner's theorem); Volume~4, Chapter~4 (Weierstrass's theorem); Volume~2, Chapter~4.
\leads Schoenberg (1928) proved that the limiting distribution function exists and is continuous, and Erdős (1939) that it is
purely singular: it has no density, although it increases strictly on $[0,1]$. A numerical calculation draws it; the proportion of $n$ with
$\varphi(n)\le n/2$ is about $0.511$: the even numbers, and a few odd numbers with many small prime factors, such as
$105=3\cdot5\cdot7$, for which $\varphi(105)/105=\frac{48}{105}$.
\end{solution}

\begin{solution}{pb:visible}
\idea Each non-zero lattice point is $d$ times a visible point, with $d$ the gcd of its coordinates. So the counting function of all
points is a sum of rescaled counting functions of visible points, and Möbius inversion goes back.

\writeup (a) If $d=\gcd(a,b)>1$, the point $(a/d,b/d)$ lies strictly between $0$ and $(a,b)$. Conversely, a lattice point on the segment
is $t(a,b)$ with $0<t<1$; if $t=\frac uv$ in lowest terms, $v\mid ua$ and $v\mid ub$, so $v\mid\gcd(a,b)$, and $v>1$.

(b) Let $N(R)$ be the number of lattice points in the closed disc of radius $R$. By Volume~6, Exercise~2.4, $|N(R)-\pi R^2|\le C(R+1)$.
Writing each non-zero point as $d\cdot v$ with $v$ visible, $N(R)-1=\sum_{d\ge1}V(R/d)$, and $V(R/d)=0$ for $d>R$. By Möbius inversion
(the sum is finite),
\[
V(R)=\sum_{d\le R}\mu(d)\bigl(N(R/d)-1\bigr)=\pi R^2\sum_{d\le R}\frac{\mu(d)}{d^2}+O\Bigl(\sum_{d\le R}\Bigl(\frac Rd+1\Bigr)\Bigr).
\]
The error is $O(R\log R)$, and $\sum_{d\le R}\frac{\mu(d)}{d^2}=\frac6{\pi^2}+O(\frac1R)$. So $V(R)=\frac6\pi R^2+O(R\log R)$.

(c) Repeat (b) with the sector $\Sigma$ of angle $\theta$ and radius $R$ in place of the disc. The unit squares centred at the lattice
points of $\Sigma$ cover $\Sigma$ minus a band of width $\frac{\sqrt2}2$ along its boundary, and lie in $\Sigma$ plus such a band; the
boundary has length $2R+\theta R$, so the number of lattice points in $\Sigma$ is $\frac\theta2R^2+O(R+1)$. The Möbius inversion is
unchanged (a sector is mapped to the same sector by $v\mapsto dv$), and gives $\frac{6}{\pi^2}\frac\theta2R^2+O(R\log R)$ visible
points. Divide by $V(R)$.

\tools Volume~6, Chapter~2 (Gauss's circle problem); Volume~2, Chapter~4 (Möbius inversion).
\leads The visible points have holes of every size: by Chinese remainders there are $k\times k$ squares of invisible points for every
$k$. In higher dimension the density of visible points is $\frac1{\zeta(n)}$. The gaps between the angles of visible points have a
limiting distribution computed by Boca, Cobeli and Zaharescu (2000), which is not exponential: the angles are equidistributed but
not Poisson.
\end{solution}

\begin{solution}{pb:farey}
\idea Neighbours in the Farey sequence span a parallelogram of area $1$ in the lattice. Bézout produces the candidate neighbour,
and a short inequality shows that nothing fits in between.

\writeup (a) The fractions of $\mathcal F_N$ with denominator $k$ are $\frac hk$ with $0\le h\le k$, $\gcd(h,k)=1$: there are $\varphi(k)$ of
them for $k\ge2$ and two ($\frac01,\frac11$) for $k=1$. So $|\mathcal F_N|=1+\sum_{k\le N}\varphi(k)$, and Mertens's estimate
$\sum_{k\le N}\varphi(k)=\frac3{\pi^2}N^2+O(N\log N)$ (Volume~6, Theorem~2.6) concludes.

(b) Let $\frac hk<1$ be in $\mathcal F_N$. Since $\gcd(h,k)=1$, the solutions $(x,y)$ of $kx-hy=1$ form one family $y\equiv y_0\pmod k$,
so there is one with $N-k<y\le N$; then $\gcd(x,y)=1$ and $\frac xy=\frac hk+\frac1{ky}$, and $0<\frac xy\le1$, because $h\le k-1$ gives
$\frac xy\le\frac{k-1}k+\frac1{ky}\le\frac{k-1}k+\frac1k=1$. So $\frac xy\in\mathcal F_N$. Suppose some $\frac{h'}{k'}\in\mathcal F_N$ lies strictly
between. Then $\frac xy-\frac{h'}{k'}\ge\frac1{yk'}$ and $\frac{h'}{k'}-\frac hk\ge\frac1{kk'}$ (positive fractions with these
denominators), so
\[
\frac1{ky}=\frac xy-\frac hk\ge\frac1{yk'}+\frac1{kk'}=\frac{k+y}{kyk'}>\frac N{kyk'}\ge\frac1{ky},
\]
because $k+y>N\ge k'$. This is absurd, so $\frac xy$ is the successor of $\frac hk$: $h'=x$, $k'=y$, $h'k-hk'=1$, and $k+k'>N$.

(c) The gap between consecutive fractions is $\frac{h'}{k'}-\frac hk=\frac1{kk'}$, and the gaps add up to $1-0$. Since $k,k'\le N$, $k\ne k'$
unless $k=k'=1$ (consecutive fractions with the same denominator $k\ge2$ would differ by $\frac1{kk}$, contradicting $h'k-hk'=1$ with
$k\mid1$), so for $N\ge2$ each gap satisfies $\frac1{kk'}\ge\frac1{N(N-1)}$, and $\frac1{kk'}\le\frac1{k\cdot(N+1-k)}\le\frac1N$ by
$k+k'\ge N+1$.

\tools Volume~2, Chapter~1 (Bézout); Volume~6, Chapter~2 (Mertens's estimate for $\sum\varphi$).
\leads Franel (1924) and Landau showed that the Riemann hypothesis is equivalent to the statement that the Farey fractions are
very evenly spread: $\sum_j|f_j-\frac jn|=O(N^{1/2+\varepsilon})$, where $f_1<\dots<f_n$ is $\mathcal F_N$. The mediant
$\frac{h+h'}{k+k'}$ of two neighbours is the first fraction to appear between them, which is how the Stern--Brocot tree of
Chapter~7 is built.
\end{solution}

\begin{solution}{pb:orderedfact}
\idea An ordered factorisation into $k$ factors $>1$ is counted by the Dirichlet series $(\zeta(s)-1)^k$, and the geometric series in
$\zeta-1$ converges exactly when $\zeta(s)<2$. The number $\rho$ where $\zeta(\rho)=2$ is the natural growth exponent, and a direct
induction shows $H(n)\le n^\rho$.

\writeup (a) For real $s>1$, $\zeta(s)-1=\sum_{n\ge2}n^{-s}$, so $(\zeta(s)-1)^k=\sum_nH_k(n)n^{-s}$, where $H_k(n)$ counts the ordered
factorisations of $n$ into exactly $k$ factors $>1$. Since $H=\sum_kH_k$ (with $H_0=\mathbf1_{n=1}$) and all terms are non-negative,
$\sum_nH(n)n^{-s}=\sum_k(\zeta(s)-1)^k$, which converges to $\frac1{2-\zeta(s)}$ when $\zeta(s)<2$, that is, when $s>\rho$ ($\zeta$ is decreasing).

(b) An ordered factorisation of $p^a$ is a composition $a=a_1+\dots+a_k$ with $a_i\ge1$, and there are $2^{a-1}$ compositions of $a$.
For $n=p_1\cdots p_k$, an ordered factorisation is an ordered partition of $\{p_1,\dots,p_k\}$ into non-empty blocks: $1,3,13,75$
for $k=1,2,3,4$ (the Fubini numbers; for $k=3$: $1$ factorisation with one factor, $6$ with two, $6$ with three).

(c) The first factor of an ordered factorisation of $n\ge2$ is a divisor $e>1$, and the rest is an ordered factorisation of $n/e$, so
$H(n)=\sum_{d\mid n,\,d<n}H(d)$. Assume $H(d)\le d^\rho$ for $d<n$ ($H(1)=1$). Then
\[
H(n)\le\sum_{d\mid n,\,d<n}d^\rho=n^\rho\sum_{e\mid n,\,e>1}e^{-\rho}<n^\rho\sum_{e\ge2}e^{-\rho}=n^\rho(\zeta(\rho)-1)=n^\rho.
\]

(d) For $\sigma>\rho$ and $n\le x$, $1\le(x/n)^\sigma$, so $\sum_{n\le x}H(n)\le x^\sigma\sum_nH(n)n^{-\sigma}=\frac{x^\sigma}{2-\zeta(\sigma)}$.

\tools Volume~6, Chapter~2 (Dirichlet series, convolution); Volume~5, Chapter~12 (the function $\zeta$ on $]1,\infty[$).
\leads Kalmár proved in 1931 that $\sum_{n\le x}H(n)\sim-\frac{x^\rho}{\rho\zeta'(\rho)}=0.3181\ldots x^\rho$, by a Tauberian theorem
applied at the simple pole of $\frac1{2-\zeta(s)}$ at $s=\rho$; a direct computation verifies it. The bound (c) is due to Chor, Lemke and
Mador (2000), and the exponent $\rho$ cannot be lowered: Hille showed in 1936 that $H(n)\ge n^{\rho-\varepsilon}$ for infinitely many $n$.
\end{solution}

\begin{solution}{pb:fracmean}
\idea $\{\frac xn\}=\frac xn-\lfloor\frac xn\rfloor$, and both sums are classical: the first is $x(\log x+\gamma)$, the second is the
divisor sum $x\log x+(2\gamma-1)x$. The difference keeps only the constants.

\writeup (a) By Volume~6, Lemma~2.5(a), $\sum_{n\le x}\frac xn=x\log x+\gamma x+O(1)$. Counting the pairs $(n,m)$ with $nm\le x$ in two ways,
$\sum_{n\le x}\lfloor\frac xn\rfloor=\sum_{m\le x}\tau(m)=x\log x+(2\gamma-1)x+O(\sqrt x)$ by Dirichlet's theorem (Volume~6, Chapter~2, via the hyperbola
method, Lemma~2.8). Subtract: $\sum_{n\le x}\{\frac xn\}=(1-\gamma)x+O(\sqrt x)$.

(b) $N\bmod n=N-n\lfloor\frac Nn\rfloor$, and $\sum_{n\le N}n\lfloor\frac Nn\rfloor=\sum_{nm\le N}n=\sum_{k\le N}\sigma(k)=\frac{\pi^2}{12}N^2+O(N\log N)$ by Volume~6,
Theorem~2.6(b). So the sum is $N^2-\frac{\pi^2}{12}N^2+O(N\log N)$.

(c) With $t=\frac nx$, $\frac1x\sum_{n\le x}\{\frac xn\}$ is a Riemann sum of $\{\frac1t\}$ on $]0,1]$, whose integral is $1-\gamma$ (Problem~\ref{pb:fracint}). Likewise
$\frac2{N^2}\sum_nn\{\frac Nn\}=\frac2{N^2}\sum(N\bmod n)$ is a Riemann sum of $2t\{\frac1t\}$, and by (b) it tends to $2-\frac{\pi^2}6$; so
$\int_0^12t\{\frac1t\}\dd t=2-\frac{\pi^2}6$, which can also be checked directly as in Problem~\ref{pb:fracint}: $2\int_1^\infty\{u\}u^{-3}\dd u=2-\zeta(2)$.
These Riemann-sum readings are heuristics, since $\{\frac1t\}$ jumps infinitely often near $0$; the error terms come from (a) and (b),
that is, from number theory.

\tools{Volume~6, Lemmas~2.5 and~2.8, Theorem~2.6; Problem~\ref{pb:fracint}.}
\leads The error in (a) is $-\Delta(x)+O(1)$, where $\Delta$ is the error in Dirichlet's divisor problem; so the question ``how regular
are the remainders of $x$ modulo $n$?'' is the divisor problem itself, open since 1849.
\end{solution}

\begin{solution}{pb:radical}
\idea The radical forgets repeated prime factors, and only numbers with repeated factors make $\frac{\operatorname{rad}(n)}n$ differ from $1$. Squarefull numbers are rare,
about $\sqrt x$ of them, so the convolution trick of Volume~6, Chapter~2, works with a power saving.

\writeup (a) $f(n)=\frac{\operatorname{rad}(n)}n$ is multiplicative with $f(p^k)=p^{1-k}$. Möbius inversion gives a multiplicative $g=f*\mu$ with $g(p^k)=f(p^k)-f(p^{k-1})$:
\[
g(p)=1-1=0,\qq g(p^k)=p^{1-k}-p^{2-k}=-(p-1)p^{1-k}\quad(k\ge2).
\]
So $g(d)=0$ unless $d$ is squarefull, and then $|g(d)|=\prod_{p^k\parallel d}(p-1)p^{1-k}\le1$, since $p-1\le p^{k-1}$ for $k\ge2$.

(b) Every squarefull $d$ can be written $d=a^2b^3$ (put one or three copies of each prime with odd exponent into $b^3$). For fixed $b$ there are at most $\sqrt{y/b^3}$ values of
$a$, so the count is at most $\sqrt y\sum_bb^{-3/2}=\zeta(\frac32)\sqrt y$.

(c) Since $\operatorname{rad}(n)=n\sum_{d\mid n}g(d)$, writing $n=dm$,
\[
\begin{aligned}\sum_{n\le x}\operatorname{rad}(n)&=\sum_{d\le x}g(d)\,d\sum_{m\le x/d}m=\sum_{d\le x}g(d)\,d\Bigl(\frac{x^2}{2d^2}+O\Bigl(\frac xd\Bigr)\Bigr)\\&=\frac{x^2}2\sum_{d\le x}\frac{g(d)}d+O\Bigl(x\sum_{d\le x}|g(d)|\Bigr).\end{aligned}
\]
By (a) and (b), $\sum_{d\le x}|g(d)|\le\zeta(\frac32)\sqrt x$, and by partial summation
\[
\sum_{d>x}\frac{|g(d)|}d\le\sum_{\substack{d>x\\ d\text{ squarefull}}}\frac1d=O(x^{-1/2}).
\]
So
$\sum_{n\le x}\operatorname{rad}(n)=\frac{x^2}2G+O(x^{3/2})$ with $G=\sum_d\frac{g(d)}d$. The series converges absolutely, so it is an Euler product:
\[
G=\prod_p\Bigl(1-(p-1)\sum_{k\ge2}\frac{p^{1-k}}{p^k}\Bigr)=\prod_p\Bigl(1-\frac{(p-1)p^{-3}}{1-p^{-2}}\Bigr)=\prod_p\Bigl(1-\frac1{p(p+1)}\Bigr)=C.
\]
The mean value of $f=\mathbf1*g$ is $\sum_d\frac{g(d)}d=C$ by Wintner's theorem (Volume~6, Theorem~2.2). Numerically, one finds $\frac2{x^2}\sum_{n\le10^7}\operatorname{rad}(n)=0.704446$, against $C=0.704442$.

\tools{Volume~6, Theorem~2.2 (Wintner) and the convolution method of Chapter~2; Euler products (Volume~5, Theorem~12.3).}
\leads Both the mean of $\frac{\operatorname{rad}(n)}n$ and the mean of $\operatorname{rad}(n)$ relative to $\frac x2$ equal $C$: the weight $n$ only rescales by $\frac12$. So a typical radical
is a fixed fraction of $n$, and numbers with a small radical, the squarefull ones, are rare. The $abc$ conjecture asserts that when $a+b=c$ with coprime $a,b$, the radical
of $abc$ can be much smaller than $c$ only in a weak sense: $c\le K_\varepsilon\operatorname{rad}(abc)^{1+\varepsilon}$.
\end{solution}

\section*{Sources and further reading}

\begin{itemize}
\item T.~M.~Apostol, \emph{Introduction to Analytic Number Theory}, Springer, 1976, Chapter~3 (averages of arithmetic functions, the
hyperbola method, squarefree integers, visible lattice points).
\item G.~H.~Hardy and E.~M.~Wright, \emph{An Introduction to the Theory of Numbers}, 6th ed., Oxford, 2008, Chapter~XVIII (the order of
magnitude of arithmetic functions, Wigert's theorem).
\item G.~Tenenbaum, \emph{Introduction to Analytic and Probabilistic Number Theory}, 3rd ed., AMS GSM~163, 2015, Chapter~I.3 (average orders) and
Part~III (mean values of arithmetic functions from the probabilistic point of view).
\item A.~Wintner, \emph{Eratosthenian Averages}, Waverly Press, Baltimore, 1943.
\item E.~C.~Titchmarsh, \emph{The Theory of the Riemann Zeta-Function}, 2nd ed., revised by D.~R.~Heath-Brown, Oxford, 1986, Chapter~XII
(the divisor problem: Vorono\"\i, Hardy, Cram\'er).
\item M.~N.~Huxley, \emph{Exponential sums and lattice points III}, Proceedings of the London Mathematical Society 87 (2003), 591--609.
\item H.~L.~Montgomery and R.~C.~Vaughan, \emph{The distribution of squarefree numbers}, in \emph{Recent Progress in Analytic Number Theory},
Vol.~1, Academic Press, 1981, 247--256.
\end{itemize}
